% BHU PYQ -- Manually transcribed & error-corrected
\documentclass[11pt,a4paper]{article}

\usepackage[a4paper,top=2.5cm,bottom=2.5cm,left=2.8cm,right=2.8cm,headheight=14pt]{geometry}
\usepackage{amsmath,amssymb}
\usepackage[T1]{fontenc}
\usepackage[utf8]{inputenc}
\usepackage{lmodern}
\usepackage{microtype}
\usepackage{fancyhdr}
\usepackage{enumitem}
\usepackage{hyperref}
\usepackage{lastpage}
\usepackage{parskip}

\hypersetup{colorlinks=true,urlcolor=black,linkcolor=black,
  pdftitle={CHB-04A: Ancillary Chemistry-II -- BHU PYQ},pdfauthor={Banaras Hindu University}}

\pagestyle{fancy}\fancyhf{}
\renewcommand{\headrulewidth}{0.4pt}\renewcommand{\footrulewidth}{0.4pt}
\fancyhead[L]{\small\textit{Banaras Hindu University}}
\fancyhead[R]{\small\textit{CHB-04A: Ancillary Chemistry-II}}
\fancyfoot[C]{\small Page \thepage\ of \pageref{LastPage}}

\newlist{parts}{enumerate}{2}
\setlist[parts,1]{label=(\alph*),leftmargin=*,itemsep=5pt,topsep=3pt}
\setlist[parts,2]{label=(\roman*),leftmargin=*,itemsep=3pt,topsep=2pt}
\newcommand{\pts}[1]{\hfill\textit{[#1]}}

\begin{document}

\begin{center}
  {\large\bfseries BANARAS HINDU UNIVERSITY}\\[2pt]
  {\normalsize B.Sc. (Hons.) Semester IV Examination, 2022-23}\\[6pt]
  \rule{0.8\linewidth}{0.5pt}\\[6pt]
  {\large\bfseries Paper No. CHB-04A: Ancillary Chemistry-II (Basic Aspects of Chemistry)}\\[4pt]
  {\normalsize Time: 3 Hours\hfill Maximum Marks: 70}
\end{center}
\rule{\linewidth}{0.4pt}
\small
\noindent\textit{Instructions: Answer five questions in total, including Question~1 which is compulsory.}
\normalsize
\bigskip

\noindent\textbf{Question 1.} Answer all parts of the following: \pts{5 $\times$ 2 = 10}
\begin{parts}
  \item ``Petroleum and natural gas are stored forms of solar energy.'' Justify this statement.
  \item Define Kranz anatomy. Draw a neat labeled structural representation of Kranz anatomy in a leaf.
  \item Discuss the different factors that affect the rate of photosynthesis in plants.
  \item Write down the chemical structures of Nylon-6,6 and Bakelite.
  \item Differentiate between a cofactor and a prosthetic group.
\end{parts}

\medskip\rule{\linewidth}{0.2pt}

\noindent\textbf{Question 2.}
\begin{parts}
  \item What is a fuel cell? Give the electrode reactions of a hydrogen-oxygen fuel cell. \pts{5}
  \item Discuss intermolecular and intramolecular hydrogen bonds with one example of each. \pts{5}
  \item What is corrosion and what are its classifications? Discuss the mechanism of rust formation on iron. How can it be prevented by cathodic protection? \pts{4}
\end{parts}

\medskip\rule{\linewidth}{0.2pt}

\noindent\textbf{Question 3.}
\begin{parts}
  \item Differentiate between isothermal and adiabatic processes. How do changes in enthalpy ($\Delta H$), entropy ($\Delta S$), and Gibbs free energy ($\Delta G$) affect the spontaneity of a thermodynamic process? \pts{5}
  \item Give two examples each of extensive and intensive thermodynamic properties. \pts{4}
  \item What is a protein? State the significance of the primary and secondary structures of a protein. \pts{5}
\end{parts}

\medskip\rule{\linewidth}{0.2pt}

\noindent\textbf{Question 4.}
\begin{parts}
  \item Define the Respiratory Quotient (RQ). Why is the RQ value for a carbohydrate substrate unity, while for substrates like fats and proteins it is less than one? \dots \pts{5}
  \item What is an enzyme? Discuss the types and roles of enzymes in living things. \pts{4}
  \item Explain the Born-Haber cycle for the formation of ionic $\text{NaCl}$. What are the factors affecting the lattice energy of ionic solids? Among $\text{NaCl}$ and $\text{MgO}$, which has higher lattice energy and why? \pts{5}
\end{parts}

\medskip\rule{\linewidth}{0.2pt}

\noindent\textbf{Question 5.}
\begin{parts}
  \item What are the different types of thermodynamic systems? Explain them briefly with suitable examples. \pts{6}
  \item How is the standard electrode potential ($E^\circ$) different from the single electrode potential ($E$)? Write the Nernst equation for a single electrode. \pts{3}
  \item What is petroleum refining? Why is it carried out? Name the different fractions obtained during refining of petroleum along with their boiling ranges and carbon atom ranges. \pts{5}
\end{parts}

\medskip\rule{\linewidth}{0.2pt}

\noindent\textbf{Question 6.}
\begin{parts}
  \item What do you mean by the vulcanization of rubber? Write down the preparation method, properties, and applications of Polyneoprene and Buna-S rubber. \pts{5}
  \item What is the significance of the octane number and cetane number in petrol and diesel engines? Draw the structures of the organic compounds having the highest and lowest values of octane number. \pts{5}
  \item Define aerobic and anaerobic respiration with suitable chemical equations. \pts{4}
\end{parts}

\medskip\rule{\linewidth}{0.2pt}

\noindent\textbf{Question 7.}
\begin{parts}
  \item Describe the biological role and structural differences between starch and cellulose. \pts{4}
  \item Describe fibrous and globular proteins and their functions. \pts{4}
  \item Differentiate between the following pairs in brief: \pts{4 $\times$ 1.5 = 6}
  \begin{parts}
    \item Respiration and Photorespiration
    \item Coordinate covalent bond and Covalent bond
    \item Photosystem I and Photosystem II in photosynthesis
    \item Extensive properties and Intensive properties
  \end{parts}
\end{parts}

\vfill
\rule{\linewidth}{0.4pt}
\begin{center}\small
  \href{https://bhu-exam.vercel.app/}{Click here to visit our website}\\[4pt]
  \href{https://me-aryan.vercel.app/}{Click here to visit our contributor}\\[4pt]
  \href{https://quashan-me.vercel.app/}{Click here to visit our contributor}
\end{center}

\end{document}
